\section{A sharp bound and local arithmetic for suffix weight folds}
\label{sec:suffix-folds}

The weighted-transport chapter of the research archive reviewed in
Section~\ref{sec:compact-coordinates} analyses a full-carrier transfer.
It moves all weights from the rightmost carrier range, including points
that survive the next truncation, and proves a sharp increase of at most
three weight runs. The maintained implementation instead moves only the
points actually removed by truncation. We prove a sharp increase of at
most two runs for this existing operation, then exploit its small receiving
interval to avoid redoing vector arithmetic on the untouched prefix.
The complete weighted certificates and the independent verifier are unchanged.

\subsection{The native two-boundary theorem}

Let the current universe be $[0,N)$, with a canonical complete partition
into $R$ constant integer-vector runs. Suppose a legal truncation retains
$[0,C)$, where $0<C<N$. Write $J$ for the interior change points,
so $R=1+|J|$, and split it into $J_- = J\cap(0,C)$ and
$J_+ = J\cap(C,N)$. A possible old change at $C$ disappears.

For a positive carrier with period $p$, the native transport sends a
removed point $x\ge C$ to
\[
 \rho(x)=b+((x-b)\bmod p),\qquad b=C-p.
\]
Every image belongs to $[b,C)$. This is iteration of the carrier's inverse
until first reaching the retained prefix: for an overlapping carrier all
intermediate inverse steps are defined, and for a disjoint carrier only
one step is needed. These are the same partial-isometry facts used in the
existing truncation proof. The preexisting prefix weights stay in place.

Extend the removed suffix weights by zero. Their discrete derivative has
support only at $C$, $N$ and $J_+$. Summation over residue classes maps
these changes into $[b,C)$; the change at $C$ maps to $b$. Extending that
residue sum by zero below $b$ can introduce a boundary at $b$, while $C$
is the endpoint of the new universe. Thus every new interior change belongs
to
\[
 J_-\ \cup\ \rho(J_+)\ \cup\ \{b,\rho(N)\},
\]
with endpoints outside $(0,C)$ discarded. At most two new boundaries
replace the vanished suffix boundaries, and the new run count satisfies
\[
 R'\le R+2.
\]
If $b=0$, or if $p$ divides $N-C$ so the two extra points coincide, the
increase is at most one. Simultaneous cancellation of vector jumps only
reduces these bounds. The reasoning applies to signed vectors without
any positivity assumption.

For a reversing carrier with left domain endpoint $a$, the removed suffix
maps through $x\mapsto a+N-1-x$ onto $[a,a+N-C)\subseteq[0,C)$.
Interior suffix changes move to $a+N-j$ for $j\in J_+$. The only extra
potential changes are the receiving interval's two endpoints, so again
$R'\le R+2$. An endpoint at $0$ or $C$ is not an interior change and
improves the bound. Other trace operations leave point weights fixed.
Static contraction takes an ordered subsequence of values, which cannot
increase the number of runs. Consequently, after $T$ native truncations,
\[
                  R_T\le R_0+2T.
\]
This sharpens the qualitative polynomial accounting in
Section~\ref{sec:weighted-components}; it does not change the orbit schedule.

The report's legal seven-point example distinguishes the two operations.
Give every point weight one and take inclusive pairings
\[
 [1,2]\longrightarrow[5,6],\quad
 [3,4]\longrightarrow[4,5],\quad
 [0,0]\longrightarrow[3,3].
\]
The maintained scheduler first truncates $N=7$ to $C=6$ using the first
carrier, of period four. Only point six is removed, and it maps to point
two. The native result is
\[
 (1\text{ on }[0,2)),\quad(2\text{ on }[2,3)),\quad
 (1\text{ on }[3,6)),
\]
so one run becomes three and the $+2$ bound is attained. The report's full
range transfer moves both points five and six, leaves point five at zero,
and produces four runs, attaining its $+3$ bound. Its claim that full-range
transfer is essential to a small run bound must therefore be read as
specific to that delivered analysis; our deleted-suffix operation has the
stronger bound above. We retain the native operation and its separate checker,
rather than adopting the report's shared producer/checker arithmetic.

\subsection{Encoded-state bounds}

For additive source intervals $[l_i,h_i)$ with vectors $v_i$, put
\[
 A=\sum_i(h_i-l_i)\lVert v_i\rVert_\infty.
\]
Each original point contribution moves to exactly one representative in its
orbit. Therefore every intermediate point-weight coordinate has absolute value at most
$A$, even with overlapping source intervals and signed cancellation. Jump
coordinates can be twice as large, which adds only a constant number of bits. If the
input has $m$ intervals, maximum coordinate bit length $L$, and universe
size $N$, the required coordinate bit length is
$O(1+L+\log(m+1)+\log(N+1))$. Quotient/remainder transfers perform no loop
proportional to a quotient or to the number of represented points.

The initial sweep has $R_0\le2m+1$ runs. If the input has $k$ pairings,
no trace operation increases their number. A static contraction has at most
$2k+1$ gaps, from the complement of a union of at most $2k$ domain and range
intervals. Intersecting these gaps with $R$ weight runs produces at most
$R+4k+2$ emitted blocks. With $E$ trace events, the total emitted blocks,
and hence the number of distinct final weight vectors, are bounded by
\[
             E(R_0+2E+4k+2).
\]
Multiplicities use at most $O(\log(N+1))$ bits. This is a coarse general
weighted-output bound; the sharper normal-surface type bound of the previous
section concerns a special geometric input class.

Together with the checked polynomial orbit trace, these bounds explicitly
account for intermediate partitions, integer magnitudes and encoded output.
Classical AHT already prove polynomial weighted orbit computation, with a
four-boundary transfer bound in Lemma~15 and polynomial accounting in
Theorem~16\cite{aht-orbits}. The present suffix-specific statement improves
the maintained observer's accounting. It supplies no bound on the number
of geometric hierarchy states or on complete unknot-recognition search.

\subsection{Restricting vector arithmetic to the receiving interval}

Previously, every fold rebuilt a difference map for the entire retained
prefix, adding and subtracting every retained run's vector at its endpoints.
This repeated dense $d$-coordinate arithmetic even on runs far outside the
receiving interval. The new producer preserves the same quotient/remainder
range additions for the removed suffix. It canonicalizes only those additions
on their receiving interval and merges that small partition with the existing
runs there. Untouched prefix and tail runs reuse their immutable vector
tuples. Coalescing at the receiving boundaries preserves the exact canonical
partition; no approximate or lazily unresolved weight is returned.

Let $S$ be the number of suffix runs and $P$ the number of retained runs
intersecting the receiving interval. The new fold uses $O(R)$ traversal
and copying work, $O(S\log(S+1))$ endpoint comparisons and
$O((S+P)d)$ coordinate arithmetic and equality work, including boundary
coalescing. The old prefix reconstruction paid $O(Rd)$ vector work even
when $S+P$ was small. Both retain flat canonical run lists, so neither
avoids the linear traversal cost. Binary division, multiplication, integer
comparison and endpoint bit lengths remain part of the bit model. This is
an improvement in work on sparse affected regions, not a new scheduling rule.

Only \code{weighted\_orbits.py} changes among pre-existing runtime modules.
The checker continues its separate residue-anchor/discrete-derivative
transport and recomputes the complete histogram. Schemas, trace operations,
normal-coordinate decoding, multiplicities and statistics are unchanged.
Tests cover 2,000 signed translation folds and 2,000 signed reflection folds
against direct finite pushforwards, including the exact $+2$ run bound.
The seven-point example is tested on an actual maintained trace. Other checks
exercise zero weights, immutable prefix reuse, late cancellation and replay
with the producer's new overlay disabled.

\subsection{Native audit and performance evidence}

The driver \code{weighted\_research/suffix\_folds.py} compares against
\code{e052c7ef07f3}, the compact-coordinate implementation from the previous
section. It records 444 current source/fixture hashes and checks them before
and after each run. The baseline package is loaded from Git in a separate
namespace; only its weighted producer differs among pre-existing runtime
modules. All independent checker files are byte-for-byte unchanged.

All 1,128 maintained tests pass in 214.506 seconds. The weighted audit passes
10,000 finite systems, split evenly between the classical and Fine--Wilf
periodic rules, plus exact huge binary examples. The independent checker
passes 2,500 literal cases, 366 mutation checks, producer-disabled replay,
and a 20,002-bit endpoint case. The normal audit passes 5,100 comparisons
and certificate replays on the 1,275-vector, 45-triangulation Regina corpus.
Additional baseline comparisons give identical complete answers, including
certificates and statistics, for 2,550 normal queries and 1,000 abstract
systems. The successful audit takes 48.158 seconds while overlapping the
suite, so its duration is not a performance measurement. The outer record's
baseline and source hashes identify this run; a nested kernel record retains
its historical baseline metadata.

The first audit-driver attempt used dictionary pairings for the isolated old
producer, whose public API requires its own \code{IntervalPairing} class.
That harness error was corrected before the successful pinned audit; the
initial log is retained separately. It was not a mathematical or runtime
regression. The read-only check of \code{openai/math} found its head unchanged
at \code{fd4aeeb2ee4f}; no new upstream mathematics was imported in this step.

The complete-query benchmark runs serially after the audit, suite and first
publication. It takes 108.073 seconds and retains all 160 measured calls and
32 warm-ups. Eight supplied normal surfaces/modes use five shuffled rounds
with an A/A copy of both engines. Every timed call includes fresh geometry,
weight construction, orbit discovery, certificate production and independent
verification. Input-fixture construction, serialization and equality checks
are outside the timer. Full answers, including every certificate field and
statistic, agree on every call.

\begin{center}\small
\begin{tabular}{lrrrrr}
\toprule
Query & Old ms & New ms & Old/new & Old A/A & New A/A \\
\midrule
meridian-1 & 0.976 & 0.939 & 1.021 & 0.990 & 0.981 \\
meridian-16 & 21.757 & 21.124 & 1.038 & 1.017 & 1.012 \\
meridian-64 & 449.308 & 409.737 & 1.094 & 0.988 & 0.986 \\
meridian-128 & 2773.559 & 2462.575 & 1.126 & 1.000 & 1.001 \\
triangle-only & 4.151 & 4.172 & 1.000 & 0.999 & 1.007 \\
binary-1024 & 17.024 & 16.230 & 1.050 & 1.000 & 0.996 \\
disk-16 & 17.191 & 17.405 & 1.019 & 0.997 & 1.010 \\
disk-128 & 1213.349 & 1206.785 & 1.008 & 1.000 & 1.000 \\
\bottomrule
\end{tabular}
\end{center}


On the 128-tetrahedron full-coordinate query, median total time changes from
2.774 to 2.463 seconds, with a within-round median ratio of $1.126$.
Production medians are 1.497 and 1.225 seconds. Verification uses identical
code and proof; variations in its measured time are not an algorithmic gain.
At 64 tetrahedra the total paired ratio is $1.094$. Triangle-only queries
are unchanged ($1.000$ after rounding), while the large binary-coordinate
fixture gives $1.050$. Three-weight disk queries give only $1.019$ and
$1.008$; these small differences should not be advertised as reliable gains.
Old A/A ratios range from $0.988$ to $1.017$ and new A/A ratios from
$0.981$ to $1.012$.

This is an incremental comparison against the compact implementation, not
against the earlier dense-coordinate version. Its gains should not be
multiplied with ratios from another run to claim a measured combined speedup.
The improvement is local to weight production, leaves automatic recognition
coverage unchanged, and does not establish a general quasipolynomial bound.
All samples and deduplicated complete certificates are retained in
\code{data/suffix-fold-benchmark.json}; the table generator is
\code{data/suffix\_fold\_tables.py}.

